% GENERATED FILE. Edit canonical lesson frontmatter, then run npm run generate:books.
% canonical-entries: 65

\chapter*{Glossary}
\addcontentsline{toc}{chapter}{Glossary}
\markboth{Glossary}{Glossary}

This glossary collects every word and phrase taught in the book. Pronunciation appears when it differs from the written form; chapter numbers show where each entry is introduced.

\noindent\begin{minipage}[t]{\linewidth}
\raggedright
\textbf{\bn{আসা}}\enspace\emph{āsā}\par
\small to come — the verb that never asks whether a man or a woman is speaking\par
\footnotesize Introduced in Chapter 7.
\end{minipage}\par\medskip

\noindent\begin{minipage}[t]{\linewidth}
\raggedright
\textbf{\bn{ভাবা}}\enspace\emph{bhābā}\par
\small to think — and the machinery that made a \textquotedblleft{}become\textquotedblright{} verb into a \textquotedblleft{}think\textquotedblright{} verb\par
\footnotesize Introduced in Chapter 8.
\end{minipage}\par\medskip

\noindent\begin{minipage}[t]{\linewidth}
\raggedright
\textbf{\bn{ভাই}}\enspace\emph{bhāi}\par
\small brother — not borrowed, not a cousin, but the same word as English \textquotedblleft{}brother,\textquotedblright{} worn down by six thousand years and two different roads\par
\footnotesize Introduced in Chapter 11.
\end{minipage}\par\medskip

\noindent\begin{minipage}[t]{\linewidth}
\raggedright
\textbf{\bn{ভালো} \bn{লাগা}}\enspace\emph{bhālo lāgā}\par
\small to like — \textquotedblleft{}good sticks to me\textquotedblright{} — and beside it \bn{ভালোবাসা}, to love\par
\footnotesize Introduced in Chapter 9.
\end{minipage}\par\medskip

\noindent\begin{minipage}[t]{\linewidth}
\raggedright
\textbf{\bn{ভাত}}\enspace\emph{bhāt}\par
\small cooked rice — the meal itself, and a Sanskrit root for \textquotedblleft{}share\textquotedblright{} that will resurface two lessons from now in a very different word\par
\footnotesize Introduced in Chapter 10.
\end{minipage}\par\medskip

\noindent\begin{minipage}[t]{\linewidth}
\raggedright
\textbf{\bn{বোঝা}}\enspace\emph{bojhā}\par
\small to understand — the third knowing, and the harmony that this time moves the spelling\par
\footnotesize Introduced in Chapter 8.
\end{minipage}\par\medskip

\noindent\begin{minipage}[t]{\linewidth}
\raggedright
\textbf{\bn{বোন}}\enspace\emph{bon}\par
\small sister — not built on the root behind English \textquotedblleft{}sister\textquotedblright{} at all, but on the very root that named this book's rice two lessons ago\par
\footnotesize Introduced in Chapter 11.
\end{minipage}\par\medskip

\noindent\begin{minipage}[t]{\linewidth}
\raggedright
\textbf{\bn{বন্ধু}}\enspace\emph{bôndhu}\par
\small friend — a word built openly on \textquotedblleft{}to bind,\textquotedblright{} and the rare case where Bengali and English share both the word and the picture behind it\par
\footnotesize Introduced in Chapter 11.
\end{minipage}\par\medskip

\noindent\begin{minipage}[t]{\linewidth}
\raggedright
\textbf{\bn{চা}}\enspace\emph{chā}\par
\small tea — a loanword that crossed from China into Bengali overland, and the drink this chapter's polite offers run on\par
\footnotesize Introduced in Chapter 10.
\end{minipage}\par\medskip

\noindent\begin{minipage}[t]{\linewidth}
\raggedright
\textbf{\bn{চোখ}}\enspace\emph{chokh}\par
\small eye — built on the same Sanskrit root as \textquotedblleft{}to see,\textquotedblright{} worn down through three centuries of Bengali spelling before it looked like this\par
\footnotesize Introduced in Chapter 12.
\end{minipage}\par\medskip

\noindent\begin{minipage}[t]{\linewidth}
\raggedright
\textbf{\bn{চশমা}}\enspace\emph{choshma}\par
\small glasses — built on the very Persian word the eye lesson once named as its only cousin\par
\footnotesize Introduced in Chapter 15.
\end{minipage}\par\medskip

\noindent\begin{minipage}[t]{\linewidth}
\raggedright
\textbf{\bn{দয়া} \bn{করে}}\enspace\emph{doya kore}\par
\small please — literally \textquotedblleft{}having done a kindness,\textquotedblright{} and the noun-plus-\bn{করা} pattern's third demonstration\par
\footnotesize Introduced in Chapter 13.
\end{minipage}\par\medskip

\noindent\begin{minipage}[t]{\linewidth}
\raggedright
\textbf{\bn{দুধ}}\enspace\emph{dudh}\par
\small milk — a word whose spelling could ask for a trailing vowel but doesn't, and a root whose one secure English cousin is not the word it looks like\par
\footnotesize Introduced in Chapter 10.
\end{minipage}\par\medskip

\noindent\begin{minipage}[t]{\linewidth}
\raggedright
\textbf{\bn{দুঃখিত}}\enspace\emph{dukkhito}\par
\small sorry — built on a word two dictionaries explain two different ways, and a prefix with a real Greek cousin\par
\footnotesize Introduced in Chapter 13.
\end{minipage}\par\medskip

\noindent\begin{minipage}[t]{\linewidth}
\raggedright
\textbf{\bn{দেখা}}\enspace\emph{dækhā}\par
\small to see, to look — and the vowel that changes while the spelling stands still\par
\footnotesize Introduced in Chapter 7.
\end{minipage}\par\medskip

\noindent\begin{minipage}[t]{\linewidth}
\raggedright
\textbf{\bn{এক} \bn{দুই} \bn{তিন} \bn{চার} \bn{পাঁচ}}\enspace\emph{ek dui tin chār pā̃ch}\par
\small one to five — and the \textquotedblleft{}two\textquotedblright{} that kept its Sanskrit vowel\par
\footnotesize Introduced in Chapter 6.
\end{minipage}\par\medskip

\noindent\begin{minipage}[t]{\linewidth}
\raggedright
\textbf{\bn{হওয়া}}\enspace\emph{hôwā}\par
\small to be, to become — the second of Bengali's two be-verbs, and the one that does what \bn{আছে} cannot\par
\footnotesize Introduced in Chapter 7.
\end{minipage}\par\medskip

\noindent\begin{minipage}[t]{\linewidth}
\raggedright
\textbf{\bn{হৃদয়}}\enspace\emph{hridoy}\par
\small heart — a word kept whole from Sanskrit, and the oldest, most secure English cousin this book has shown you\par
\footnotesize Introduced in Chapter 12.
\end{minipage}\par\medskip

\noindent\begin{minipage}[t]{\linewidth}
\raggedright
\textbf{\bn{জামা}}\enspace\emph{jama}\par
\small shirt — a second Persian loan for clothing, the same word behind Hindi's own jaamaa\par
\footnotesize Introduced in Chapter 15.
\end{minipage}\par\medskip

\noindent\begin{minipage}[t]{\linewidth}
\raggedright
\textbf{\bn{জানা}}\enspace\emph{jānā}\par
\small to know a fact — beside \bn{চেনা}, the second knowing that English merged away\par
\footnotesize Introduced in Chapter 7.
\end{minipage}\par\medskip

\noindent\begin{minipage}[t]{\linewidth}
\raggedright
\textbf{\bn{যাওয়া}}\enspace\emph{jāwā}\par
\small to go — and the verb ending that carries how much respect you are paying\par
\footnotesize Introduced in Chapter 7.
\end{minipage}\par\medskip

\noindent\begin{minipage}[t]{\linewidth}
\raggedright
\textbf{\bn{জিজ্ঞাসা} \bn{করা}}\enspace\emph{jijñāsā kôrā}\par
\small to ask — literally \textquotedblleft{}to do a wanting-to-know,\textquotedblright{} and the door into a thousand verbs\par
\footnotesize Introduced in Chapter 9.
\end{minipage}\par\medskip

\noindent\begin{minipage}[t]{\linewidth}
\raggedright
\textbf{\bn{জল}}\enspace\emph{jôl}\par
\small water — the word Chapter 7 already used to say \textquotedblleft{}drink water,\textquotedblright{} now taught on its own, beside the Bangladesh-register \bn{পানি}\par
\footnotesize Introduced in Chapter 10.
\end{minipage}\par\medskip

\noindent\begin{minipage}[t]{\linewidth}
\raggedright
\textbf{\bn{কালো}}\enspace\emph{kalo}\par
\small black — the same word Sanskrit uses for time, and for death\par
\footnotesize Introduced in Chapter 14.
\end{minipage}\par\medskip

\noindent\begin{minipage}[t]{\linewidth}
\raggedright
\textbf{\bn{কাপড়}}\enspace\emph{kapor}\par
\small cloth, clothing — a word Sanskrit's own dictionaries call homegrown, not inherited, and the flapped letter returns\par
\footnotesize Introduced in Chapter 15.
\end{minipage}\par\medskip

\noindent\begin{minipage}[t]{\linewidth}
\raggedright
\textbf{\bn{খাওয়া}}\enspace\emph{khāwā}\par
\small to eat — and, in everyday Bengali, to drink as well\par
\footnotesize Introduced in Chapter 7.
\end{minipage}\par\medskip

\noindent\begin{minipage}[t]{\linewidth}
\raggedright
\textbf{\bn{লাল}}\enspace\emph{lal}\par
\small red — a jeweler's word before it was a color word, and the identical Persian loan Hindi already teaches\par
\footnotesize Introduced in Chapter 14.
\end{minipage}\par\medskip

\noindent\begin{minipage}[t]{\linewidth}
\raggedright
\textbf{\bn{লেখা}}\enspace\emph{lekhā}\par
\small to write — and the dictionary form that is quietly a noun\par
\footnotesize Introduced in Chapter 8.
\end{minipage}\par\medskip

\noindent\begin{minipage}[t]{\linewidth}
\raggedright
\textbf{\bn{মাফ} \bn{করবেন}}\enspace\emph{maf korben}\par
\small excuse me — the same Persian loan Hindi's own \textquotedblleft{}sorry\textquotedblright{} phrase uses, softened here by the future tense rather than a command\par
\footnotesize Introduced in Chapter 13.
\end{minipage}\par\medskip

\noindent\begin{minipage}[t]{\linewidth}
\raggedright
\textbf{\bn{মুখ}}\enspace\emph{mukh}\par
\small mouth, face — a word so old that scholars still argue over which language family it belongs to\par
\footnotesize Introduced in Chapter 12.
\end{minipage}\par\medskip

\noindent\begin{minipage}[t]{\linewidth}
\raggedright
\textbf{\bn{নাক}}\enspace\emph{nāk}\par
\small nose — the direct cousin of English \textquotedblleft{}nose,\textquotedblright{} worn down by Bengali's own sound changes rather than Latin's or English's\par
\footnotesize Introduced in Chapter 12.
\end{minipage}\par\medskip

\noindent\begin{minipage}[t]{\linewidth}
\raggedright
\textbf{\bn{নেওয়া}}\enspace\emph{neowā}\par
\small to take — and the verb that finishes other verbs\par
\footnotesize Introduced in Chapter 9.
\end{minipage}\par\medskip

\noindent\begin{minipage}[t]{\linewidth}
\raggedright
\textbf{\bn{নীল}}\enspace\emph{nil}\par
\small blue — Sanskrit's word for indigo, one form for every noun it describes\par
\footnotesize Introduced in Chapter 14.
\end{minipage}\par\medskip

\noindent\begin{minipage}[t]{\linewidth}
\raggedright
\textbf{\bn{পড়া}}\enspace\emph{pôṛā}\par
\small to read — and the flapped letter Sanskrit never wrote\par
\footnotesize Introduced in Chapter 8.
\end{minipage}\par\medskip

\noindent\begin{minipage}[t]{\linewidth}
\raggedright
\textbf{\bn{পরিবার}}\enspace\emph{pôribār}\par
\small family — literally \textquotedblleft{}what surrounds you,\textquotedblright{} a second word Bengali kept whole from Sanskrit rather than wearing down\par
\footnotesize Introduced in Chapter 11.
\end{minipage}\par\medskip

\noindent\begin{minipage}[t]{\linewidth}
\raggedright
\textbf{\bn{সাহায্য} \bn{করা}}\enspace\emph{sāhājjo kôrā}\par
\small to help — \textquotedblleft{}to go together with,\textquotedblright{} and the word that finally shows you the inherent vowel\par
\footnotesize Introduced in Chapter 9.
\end{minipage}\par\medskip

\noindent\begin{minipage}[t]{\linewidth}
\raggedright
\textbf{\bn{সাদা}}\enspace\emph{shada}\par
\small white — a tadbhava worn down from Sanskrit, and the one color where Bengali and Hindi genuinely parted ways\par
\footnotesize Introduced in Chapter 14.
\end{minipage}\par\medskip

\noindent\begin{minipage}[t]{\linewidth}
\raggedright
\textbf{\bn{শাড়ি}}\enspace\emph{shari}\par
\small sari — a Sanskrit word for a strip of cloth that walked, almost unchanged, straight into English\par
\footnotesize Introduced in Chapter 15.
\end{minipage}\par\medskip

\noindent\begin{minipage}[t]{\linewidth}
\raggedright
\textbf{\bn{স্বাগতম}}\enspace\emph{shbagotom}\par
\small welcome — \textquotedblleft{}a good coming,\textquotedblright{} built on the same \textquotedblleft{}hither\textquotedblright{} prefix already met, and a verb whose PIE root is unusually secure\par
\footnotesize Introduced in Chapter 13.
\end{minipage}\par\medskip

\noindent\begin{minipage}[t]{\linewidth}
\raggedright
\textbf{\bn{সবুজ}}\enspace\emph{shobuj}\par
\small green — a second Persian loan, and the chapter's payoff on all five colors\par
\footnotesize Introduced in Chapter 14.
\end{minipage}\par\medskip

\noindent\begin{minipage}[t]{\linewidth}
\raggedright
\textbf{\bn{আচ্ছা}}\par
\small okay / alright / I see (āchchhā)\par
\footnotesize Introduced in Chapter 1.
\end{minipage}\par\medskip

\noindent\begin{minipage}[t]{\linewidth}
\raggedright
\textbf{\bn{আবার}}\par
\small again\par
\footnotesize Introduced in Chapter 4.
\end{minipage}\par\medskip

\noindent\begin{minipage}[t]{\linewidth}
\raggedright
\textbf{\bn{আবার} \bn{দেখা} \bn{হবে}}\par
\small we'll meet again (goodbye)\par
\footnotesize Introduced in Chapter 4.
\end{minipage}\par\medskip

\noindent\begin{minipage}[t]{\linewidth}
\raggedright
\textbf{\bn{আমি}}\par
\small I\par
\footnotesize Introduced in Chapter 3.
\end{minipage}\par\medskip

\noindent\begin{minipage}[t]{\linewidth}
\raggedright
\textbf{\bn{আমি} \bn{বাংলা} \bn{বলি}}\par
\small I speak Bengali\par
\footnotesize Introduced in Chapter 5.
\end{minipage}\par\medskip

\noindent\begin{minipage}[t]{\linewidth}
\raggedright
\textbf{\bn{আমার}}\par
\small my\par
\footnotesize Introduced in Chapter 2.
\end{minipage}\par\medskip

\noindent\begin{minipage}[t]{\linewidth}
\raggedright
\textbf{\bn{আমার} \bn{নাম} …}\par
\small my name is… (with no \textquotedblleft{}is\textquotedblright{})\par
\footnotesize Introduced in Chapter 2.
\end{minipage}\par\medskip

\noindent\begin{minipage}[t]{\linewidth}
\raggedright
\textbf{\bn{আলাপ} \bn{করে} \bn{ভালো} \bn{লাগলো}}\par
\small pleased to meet you (lit. \textquotedblleft{}having made acquaintance, {[}it{]} felt good\textquotedblright{})\par
\footnotesize Introduced in Chapter 2.
\end{minipage}\par\medskip

\noindent\begin{minipage}[t]{\linewidth}
\raggedright
\textbf{\bn{আসি}}\par
\small goodbye, lit. \textquotedblleft{}I'll come {[}again{]}\textquotedblright{} (āshi)\par
\footnotesize Introduced in Chapter 1.
\end{minipage}\par\medskip

\noindent\begin{minipage}[t]{\linewidth}
\raggedright
\textbf{\bn{কি}}\par
\small what\par
\footnotesize Introduced in Chapter 2.
\end{minipage}\par\medskip

\noindent\begin{minipage}[t]{\linewidth}
\raggedright
\textbf{\bn{কাজ} \bn{করা}}\par
\small to work (lit. \textquotedblleft{}work-do\textquotedblright{})\par
\footnotesize Introduced in Chapter 5.
\end{minipage}\par\medskip

\noindent\begin{minipage}[t]{\linewidth}
\raggedright
\textbf{\bn{কোনো} \bn{ব্যাপার} \bn{না}}\par
\small it's no matter / no problem / you're welcome\par
\footnotesize Introduced in Chapter 3.
\end{minipage}\par\medskip

\noindent\begin{minipage}[t]{\linewidth}
\raggedright
\textbf{\bn{কেমন}}\par
\small how\par
\footnotesize Introduced in Chapter 3.
\end{minipage}\par\medskip

\noindent\begin{minipage}[t]{\linewidth}
\raggedright
\textbf{\bn{কাল} \bn{দেখা} \bn{হবে}}\par
\small see you tomorrow (and the \bn{কাল} puzzle)\par
\footnotesize Introduced in Chapter 4.
\end{minipage}\par\medskip

\noindent\begin{minipage}[t]{\linewidth}
\raggedright
\textbf{\bn{তুমি} / \bn{আপনি}}\par
\small you (familiar / respectful; also \bn{তুই} intimate)\par
\footnotesize Introduced in Chapter 2.
\end{minipage}\par\medskip

\noindent\begin{minipage}[t]{\linewidth}
\raggedright
\textbf{\bn{তুমি} \bn{কেমন} \bn{আছো}?}\par
\small how are you? (familiar)\par
\footnotesize Introduced in Chapter 3.
\end{minipage}\par\medskip

\noindent\begin{minipage}[t]{\linewidth}
\raggedright
\textbf{\bn{তোমার} \bn{নাম} \bn{কি}?}\par
\small what's your name?\par
\footnotesize Introduced in Chapter 2.
\end{minipage}\par\medskip

\noindent\begin{minipage}[t]{\linewidth}
\raggedright
\textbf{\bn{থাকা}}\par
\small to live, to stay, to remain\par
\footnotesize Introduced in Chapter 5.
\end{minipage}\par\medskip

\noindent\begin{minipage}[t]{\linewidth}
\raggedright
\textbf{\bn{দেখা} \bn{হবে}}\par
\small (we) will meet (lit. \textquotedblleft{}seeing will happen\textquotedblright{})\par
\footnotesize Introduced in Chapter 4.
\end{minipage}\par\medskip

\noindent\begin{minipage}[t]{\linewidth}
\raggedright
\textbf{\bn{ধন্যবাদ}}\par
\small thank you (dhônyobad)\par
\footnotesize Introduced in Chapter 1.
\end{minipage}\par\medskip

\noindent\begin{minipage}[t]{\linewidth}
\raggedright
\textbf{\bn{নাম}}\par
\small name\par
\footnotesize Introduced in Chapter 2.
\end{minipage}\par\medskip

\noindent\begin{minipage}[t]{\linewidth}
\raggedright
\textbf{\bn{নমস্কার}}\par
\small hello / goodbye (nômoshkar)\par
\footnotesize Introduced in Chapter 1.
\end{minipage}\par\medskip

\noindent\begin{minipage}[t]{\linewidth}
\raggedright
\textbf{\bn{বলা}}\par
\small to speak, to say\par
\footnotesize Introduced in Chapter 5.
\end{minipage}\par\medskip

\noindent\begin{minipage}[t]{\linewidth}
\raggedright
\textbf{\bn{ভালো}}\par
\small well, good — and the reply \textquotedblleft{}\bn{আমি} \bn{ভালো} \bn{আছি}\textquotedblright{}\par
\footnotesize Introduced in Chapter 3.
\end{minipage}\par\medskip

\noindent\begin{minipage}[t]{\linewidth}
\raggedright
\textbf{\bn{হ্যাঁ} / \bn{না}}\par
\small yes / no (hyã / nā)\par
\footnotesize Introduced in Chapter 1.
\end{minipage}\par\medskip
